% GENERATED FILE. Edit canonical lessons, then run npm run generate:books.
% canonical-source-hash: fnv1a64:de931f8922f515ab
% canonical-lessons: TE-C243-pariskarincu, TE-C243-taggincu, TE-C243-pencu, TE-C243-teliyajeyu, TE-C243-badili-ceyu

\chapter{To solve, To reduce, To raise, To inform, To transfer}
\label{ch:te-to-solve-to-reduce-to-raise-to-inform-to-transfer}
\hlchaptermodality{243}

\begin{chapteropening}
\textbf{By the end of this chapter:} \emph{I can say to solve, to reduce, to raise, to inform and to transfer.}

Complete the last lesson of chapter 243: I can say to solve, to reduce, to raise, to inform and to transfer.
\end{chapteropening}

\section[\texorpdfstring{\te{పరిష్కరించు} (pariṣkariñcu)}{పరిష్కరించు (pariṣkariñcu)}]{\te{పరిష్కరించు} (pariṣkariñcu) — to solve, to settle}

\label{lesson:TE-C243-pariskarincu}



\begin{quote}
Before the new one: say the Telugu for to begin, then the Telugu for to iron.
\end{quote}

\subsection*{You'll want to know: \te{పరిష్కరించు}}
\textbf{\te{పరిష్కరించు}} — \emph{pariṣkariñcu} — \textquotedblleft{}to solve, to settle\textquotedblright{}.

\begin{grammarlens}[title={What you've built}]
One more word for this chapter.
\end{grammarlens}

\subsection*{Your turn}
\begin{itemize}
\raggedright
  \item \emph{Say it:} \emph{pariṣkariñcu}
  \item \emph{Say it:} \emph{pariṣkariñcu}, once more
  \item \emph{Say it:} say \emph{pariṣkariñcu}
  \item \emph{Recall:} say the Telugu for to begin, then the Telugu for to iron, then say \emph{pariṣkariñcu} again
\end{itemize}

\begin{tcolorbox}[breakable,colback=teal!4,colframe=teal!35!black,title={Before you move on}]
What does \te{పరిష్కరించు} mean? (\textquotedblleft{}To solve, to settle\textquotedblright{}.) Say it once more.
\end{tcolorbox}
\section[\texorpdfstring{\te{తగ్గించు} (taggiñcu)}{తగ్గించు (taggiñcu)}]{\te{తగ్గించు} (taggiñcu) — to reduce, to lower}

\label{lesson:TE-C243-taggincu}



\begin{quote}
Before the new one: say the Telugu for to iron, then the Telugu for to solve.
\end{quote}

\subsection*{You'll want to know: \te{తగ్గించు}}
\textbf{\te{తగ్గించు}} — \emph{taggiñcu} — \textquotedblleft{}to reduce, to lower\textquotedblright{}.

\begin{grammarlens}[title={What you've built}]
One more word for this chapter.
\end{grammarlens}

\subsection*{Your turn}
\begin{itemize}
\raggedright
  \item \emph{Say it:} \emph{taggiñcu}
  \item \emph{Say it:} \emph{taggiñcu}, once more
  \item \emph{Say it:} say \emph{taggiñcu}
  \item \emph{Recall:} say the Telugu for to iron, then the Telugu for to solve, then say \emph{taggiñcu} again
\end{itemize}

\begin{tcolorbox}[breakable,colback=teal!4,colframe=teal!35!black,title={Before you move on}]
What does \te{తగ్గించు} mean? (\textquotedblleft{}To reduce, to lower\textquotedblright{}.) Say it once more.
\end{tcolorbox}
\section[\texorpdfstring{\te{పెంచు} (peñcu)}{పెంచు (peñcu)}]{\te{పెంచు} (peñcu) — to raise, to increase}

\label{lesson:TE-C243-pencu}



\begin{quote}
Before the new one: say the Telugu for to solve, then the Telugu for to reduce.
\end{quote}

\subsection*{You'll want to know: \te{పెంచు}}
\textbf{\te{పెంచు}} — \emph{peñcu} — \textquotedblleft{}to raise, to increase\textquotedblright{}.

\begin{grammarlens}[title={What you've built}]
One more word for this chapter.
\end{grammarlens}

\subsection*{Your turn}
\begin{itemize}
\raggedright
  \item \emph{Say it:} \emph{peñcu}
  \item \emph{Say it:} \emph{peñcu}, once more
  \item \emph{Say it:} say \emph{peñcu}
  \item \emph{Recall:} say the Telugu for to solve, then the Telugu for to reduce, then say \emph{peñcu} again
\end{itemize}

\begin{tcolorbox}[breakable,colback=teal!4,colframe=teal!35!black,title={Before you move on}]
What does \te{పెంచు} mean? (\textquotedblleft{}To raise, to increase\textquotedblright{}.) Say it once more.
\end{tcolorbox}
\section[\texorpdfstring{\te{తెలియజేయు} (teliyajēyu)}{తెలియజేయు (teliyajēyu)}]{\te{తెలియజేయు} (teliyajēyu) — to inform, to let know}

\label{lesson:TE-C243-teliyajeyu}



\begin{quote}
Before the new one: say the Telugu for to reduce, then the Telugu for to raise.
\end{quote}

\subsection*{You'll want to know: \te{తెలియజేయు}}
\textbf{\te{తెలియజేయు}} — \emph{teliyajēyu} — \textquotedblleft{}to inform, to let know\textquotedblright{}.

\begin{grammarlens}[title={What you've built}]
One more word for this chapter.
\end{grammarlens}

\subsection*{Your turn}
\begin{itemize}
\raggedright
  \item \emph{Say it:} \emph{teliyajēyu}
  \item \emph{Say it:} \emph{teliyajēyu}, once more
  \item \emph{Say it:} say \emph{teliyajēyu}
  \item \emph{Recall:} say the Telugu for to reduce, then the Telugu for to raise, then say \emph{teliyajēyu} again
\end{itemize}

\begin{tcolorbox}[breakable,colback=teal!4,colframe=teal!35!black,title={Before you move on}]
What does \te{తెలియజేయు} mean? (\textquotedblleft{}To inform, to let know\textquotedblright{}.) Say it once more.
\end{tcolorbox}
\section[\texorpdfstring{\te{బదిలీ} \te{చేయు} (badilī cēyu)}{బదిలీ చేయు (badilī cēyu)}]{\te{బదిలీ} \te{చేయు} (badilī cēyu) — to transfer}

\label{lesson:TE-C243-badili-ceyu}



\begin{quote}
Before the new one: say the Telugu for to raise, then the Telugu for to inform.
\end{quote}

\subsection*{You'll want to know: \te{బదిలీ} \te{చేయు}}
\textbf{\te{బదిలీ} \te{చేయు}} — \emph{badilī cēyu} — \textquotedblleft{}to transfer\textquotedblright{}.

\begin{grammarlens}[title={What you've built}]
One more word for this chapter.
\end{grammarlens}

\subsection*{Your turn}
\begin{itemize}
\raggedright
  \item \emph{Say it:} \emph{badilī cēyu}
  \item \emph{Say it:} \emph{badilī cēyu}, once more
  \item \emph{Say it:} say \emph{badilī cēyu}
  \item \emph{Recall:} say the Telugu for to raise, then the Telugu for to inform, then say \emph{badilī cēyu} again
\end{itemize}

\begin{tcolorbox}[breakable,colback=teal!4,colframe=teal!35!black,title={Before you move on}]
What does \te{బదిలీ} \te{చేయు} mean? (\textquotedblleft{}To transfer\textquotedblright{}.) Say it once more.
\end{tcolorbox}
